\chapter[\emj{📏}~Propiedades de Árboles (Altura, Diámetro, Balance)]{\emj{📏}~Propiedades de Árboles (Altura, Diámetro, Balance)}

\begin{dsaquees}

Medir un árbol 🌲 preguntando cosas desde ABAJO hacia arriba. Casi todo
en un árbol se resuelve con un DFS que devuelve un dato de cada rama y lo
combina en el padre:

\begin{itemize}
\item Altura: ¿qué tan lejos está la hoja más profunda?
\item Diámetro: ¿cuál es el camino más largo entre dos hojas cualquiera?
\item ¿Está balanceado?: ¿ninguna rama es mucho más larga que su hermana?
\end{itemize}

La idea central es la \textbf{recursión post-orden}: primero resuelves
los hijos, y con esos resultados decides el del padre.

\end{dsaquees}

\begin{dsaotraforma}

Un patrón resuelve la mayoría de preguntas sobre árboles: un DFS
recursivo que retorna un valor por subárbol y lo combina en el nodo
padre.

\begin{itemize}
\item \textbf{Altura/Profundidad}: $1 + \max(h_{izq}, h_{der})$, donde
      $h$ es la altura de cada hijo.
\item \textbf{Diámetro}: el camino más largo entre dos nodos. En cada
      nodo, \verb|altura_izq + altura_der| es el diámetro que pasa por
      él; el resultado es el máximo sobre todos los nodos.
\item \textbf{Balanceado}: en cada nodo, la diferencia de alturas de sus
      hijos es $\le 1$. Se verifica en una sola pasada devolviendo la
      altura y un flag.
\item \textbf{Path Sum}: existe una raíz-a-hoja que suma un objetivo, o
      la suma máxima de camino.
\end{itemize}

La clave es calcularlo en \textbf{una sola pasada} $O(n)$, no recalcular
alturas dentro de otro recorrido (eso sería $O(n^2)$).

\end{dsaotraforma}

\begin{dsadiagrama}
\begin{verbatim}
Árbol:
          [1]
        /     \
      [2]     [3]
     /   \
   [4]   [5]

Altura = 3 (1 -> 2 -> 4)

Diámetro (camino más largo): 4 -> 2 -> 5    ó   4 -> 2 -> 1 -> 3
  En el nodo 2: altura_izq(1) + altura_der(1) = 2 aristas
  En el nodo 1: altura_izq(2) + altura_der(1) = 3 aristas  <- máximo
Diámetro = 3 aristas.

¿Balanceado? |alt(2)=2 - alt(3)=1| = 1 <= 1  -> SÍ
\end{verbatim}
\end{dsadiagrama}

\begin{dsacomofunciona}
\begin{verbatim}
CÁLCULO BOTTOM-UP: cada nodo devuelve su altura (h) hacia arriba

              (1) h=3
             /      \
       h=2 (2)      (3) h=1
          /   \
    h=1 (4)   (5) h=1

Las hojas devuelven h=1. El padre = 1 + max(hijos).
  altura(4)=1  altura(5)=1  → altura(2)=1+max(1,1)=2
  altura(3)=1                → altura(1)=1+max(2,1)=3

DIÁMETRO (se actualiza en cada nodo con hIzq + hDer, en aristas):
  en nodo 2:  hIzq(1) + hDer(1) = 2 aristas   (camino 4-2-5)
  en nodo 1:  hIzq(2) + hDer(1) = 3 aristas   (camino 4-2-1-3)  ← MÁXIMO
  Diámetro = 3

BALANCE (|hIzq - hDer| ≤ 1 en todos):
  nodo 1: |2 - 1| = 1 OK     nodo 2: |1 - 1| = 0 OK   → balanceado
\end{verbatim}
\end{dsacomofunciona}

\begin{lstlisting}
PATRÓN GENERAL (DFS post-orden que retorna un valor)
- Resuelve izquierda y derecha primero.
- Combina sus resultados para el nodo actual.
- Actualiza una respuesta global si hace falta (ej. diámetro).

PSEUDOCÓDIGO (diámetro)
──────────────────────────────────────────────────────────────

diametro = 0
función altura(nodo):
    SI nodo == NULL: return 0
    izq = altura(nodo.izq)
    der = altura(nodo.der)
    diametro = max(diametro, izq + der)   // camino que pasa por nodo
    return 1 + max(izq, der)

// llamar altura(root); la respuesta queda en 'diametro'

CÓDIGO C++
──────────────────────────────────────────────────────────────

struct Node { int val; Node *left, *right; };

// Altura -> O(n)
int height(Node* n) {
    if (!n) return 0;
    return 1 + max(height(n->left), height(n->right));
}

// Diámetro en una sola pasada -> O(n)
int diameter = 0;
int dfsDiam(Node* n) {
    if (!n) return 0;
    int l = dfsDiam(n->left), r = dfsDiam(n->right);
    diameter = max(diameter, l + r);   // aristas que cruzan por n
    return 1 + max(l, r);
}

// ¿Balanceado? -1 = no balanceado; si no, retorna la altura -> O(n)
int checkBalance(Node* n) {
    if (!n) return 0;
    int l = checkBalance(n->left);  if (l == -1) return -1;
    int r = checkBalance(n->right); if (r == -1) return -1;
    if (abs(l - r) > 1) return -1;     // desbalanceado
    return 1 + max(l, r);
}

// Path Sum raíz-a-hoja que suma 'target'
bool hasPathSum(Node* n, int target) {
    if (!n) return false;
    if (!n->left && !n->right) return target == n->val;   // hoja
    return hasPathSum(n->left,  target - n->val) ||
           hasPathSum(n->right, target - n->val);
}

COMPLEJIDAD (Big-O)
──────────────────────────────────────────────────────────────

  Propiedad     │ Tiempo  │ Espacio
  ──────────────┼─────────┼──────────
  Altura        │ O(n)    │ O(h)
  Diámetro      │ O(n)    │ O(h)
  Balanceado    │ O(n)    │ O(h)
  Path Sum      │ O(n)    │ O(h)

* h = altura (pila de recursión). Una sola pasada, no O(n²).
\end{lstlisting}

\begin{dsaentrevistas}

✅ "Maximum Depth" (104), "Diameter of Binary Tree" (543), "Balanced
   Binary Tree" (110), "Path Sum" (112): todos son el MISMO patrón DFS
   post-orden que retorna un dato por subárbol.

💡 Truco del "flag −1": para "¿está balanceado?", en vez de calcular
   altura y balance por separado (O(n²)), retorna la altura o −1 si ya
   detectaste desbalance. Una sola pasada O(n).

⚠️ Diámetro se mide en ARISTAS, no en nodos, en la mayoría de enunciados
   (LeetCode 543). Lee bien: a veces el camino no pasa por la raíz.

🔑 "Binary Tree Maximum Path Sum" (124): variante dura. En cada nodo, la
   ganancia que subes es \verb|val + max(0, izq, der)|, pero la respuesta
   global considera \verb|val + izq + der|.

\end{dsaentrevistas}

\begin{dsaresumen}

• Patrón maestro: DFS post-orden que retorna un dato por subárbol.
• Altura = 1 + max(alturas de hijos).
• Diámetro = max sobre nodos de (altura izq + altura der); en aristas.
• Balanceado: diferencia de alturas <= 1; usa flag −1 en una pasada.
• Path Sum: resta el valor del nodo al bajar; verifica en las hojas.
• Todo O(n); no recalcules alturas (evita O(n²)).
════════════════════════════════════════════════════════════════

\end{dsaresumen}
